% Generated draft: requires mathematical and source review.
\documentclass[11pt]{article}
\usepackage[margin=25mm]{geometry}
\usepackage{fontspec}
\setmainfont{Latin Modern Roman}
\usepackage{amsmath,amssymb,mathtools,unicode-math}
\setmathfont{Latin Modern Math}
\usepackage{enumitem,needspace,xurl,hyperref}
\setlength{\parindent}{0pt}
\setlength{\parskip}{5pt}
\setlength{\emergencystretch}{3em}
\newcommand{\N}{\mathbb{N}}
\newcommand{\Z}{\mathbb{Z}}
\newcommand{\Q}{\mathbb{Q}}
\newcommand{\R}{\mathbb{R}}
\newcommand{\C}{\mathbb{C}}
\newcommand{\F}{\mathbb{F}}
\newcommand{\A}{\mathbb{A}}
\newcommand{\PP}{\mathbb{P}}
\newcommand{\E}{\mathbb{E}}
\newcommand{\eps}{\varepsilon}
\newcommand{\dd}{\,\mathrm{d}}
\begin{document}
\section*{UnsolvedMath ID: 30004653}
\textbf{Source problem code:} catalogue entry
\subsection{Definitions and mathematical statement}
Let \(F\) be a field described in the supplied record as semi-global, and let \(G\) be a linear algebraic group over \(F\). The record does not specify the definition of semi-global field, the hypotheses on its underlying valued field, or additional hypotheses on \(G\). These qualifications must be fixed before a precise classification problem can be stated.

To express the local-global question conditionally, fix an index set \(I\) and field extensions \(F_i/F\) for \(i\in I\), representing the intended local fields. Neither this family nor a rule for constructing it is supplied. For the following notation, use torsors for the faithfully flat, locally finitely presented topology: a \(G\)-torsor is an \(F\)-scheme with a right \(G\)-action that becomes equivariantly isomorphic to \(G\) after a cover in that topology. Write \(H^1_{\mathrm{fppf}}(F,G)\) for the pointed set of isomorphism classes of such torsors, with distinguished element \([G]\), the trivial torsor. This makes the torsor convention explicit; the source itself does not specify a topology.

Base change defines a map of pointed sets
\[
\operatorname{res}_{(F_i)}:H^1_{\mathrm{fppf}}(F,G)\longrightarrow\prod_{i\in I}H^1_{\mathrm{fppf}}(F_i,G_{F_i}),
\qquad [P]\longmapsto([P_{F_i}])_{i\in I},
\]
where \(G_{F_i}=G\times_F F_i\) and \(P_{F_i}=P\times_F F_i\). Define the pointed set of locally trivial classes by
\[
\mathcal K_{(F_i)}(F,G)=\left\{[P]\in H^1_{\mathrm{fppf}}(F,G):[P_{F_i}]=[G_{F_i}]\text{ for every }i\in I\right\}.
\]
This is a pointed-set kernel; no group structure is assumed.

\emph{Conditional formulation of the requested problem.} Specify the intended class of semi-global fields, the permitted groups, and the local fields \((F_i)_{i\in I}\). Determine under which hypotheses the implication
\[
\left(\forall i\in I,\ P_{F_i}\simeq G_{F_i}\text{ as }G_{F_i}\text{-torsors}\right)
\Longrightarrow P\simeq G\text{ as }G\text{-torsors}
\]
holds for every \(G\)-torsor \(P\) over \(F\), equivalently when \(\mathcal K_{(F_i)}(F,G)=\{[G]\}\). Describe explicitly the obstruction to this implication. The displayed kernel merely denotes the classes witnessing failure; it does not supply the requested explicit description.

The source does not determine whether this triviality principle or the stronger requirement that \(\operatorname{res}_{(F_i)}\) be injective is intended. Nor does it specify the form of an acceptable explicit obstruction. Consequently, this is a conditional framework for the supplied question, not a uniquely determined conjecture or a claimed solution.
\subsection{Short English statement}
The question asks when a torsor over a semi-global field must be trivial if it is trivial over every intended local field, and how to describe failures explicitly. The source leaves the local fields, field and group hypotheses, and intended obstruction unspecified, so a precise unconditional target cannot be recovered. Its open status is supplied dated curation, not an independently verified mathematical conclusion.
\subsection{Sources}
\begin{itemize}
\item Supplied local problem registry record, permanent UnsolvedMath ID \(30004653\), titled \emph{Local-Global Principles for Torsors over Semi-Global Fields}; exact source problem code \texttt{catalogue entry}. The record belongs to the Oberwolfach Reports Open Problems collection and lists proposed year \(2021\); no individual proposer is supplied.
\item Supplied report citation: \emph{Homotopic and Geometric Galois Theory (2022)}, Oberwolfach Reports, pp.\ 663--744, DOI \texttt{10.4171/owr/2021/12}, \url{https://doi.org/10.4171/owr/2021/12}. The supplied literature-source label for the same link is \emph{Oberwolfach Report 12/2021}. The background describes these reports as documenting mathematical workshops and their open-problem discussions. No exact problem page or report passage is supplied, and the cited report has not been independently consulted for this draft.
\item Literature assessment checked at \texttt{2026-08-22}: this is supplied dated curation, not independently verified. The supplied status is \emph{open}. The assessment says that the extracted question identifies neither a specific obstruction nor a mathematical claim beyond an unspecified request for an explicit description, and that the cited report context does not determine a resolved target. The background expressly requests independent verification of time-sensitive status and cited claims. Its listed source checked is \emph{Oberwolfach Report 12/2021}, \url{https://doi.org/10.4171/owr/2021/12}.
\item The registry labels the statement \texttt{corrected\_verified} and reports manual semantic, typographic, and LaTeX review asserting that the clean statement is self-contained and preserves its target. These are supplied review claims; they do not resolve the missing mathematical specifications identified here.
\end{itemize}
\end{document}
